\section{Making a marked base row before choosing the target coordinates}

% BEGIN P77R1
\paragraph{P77R1. Adapting the base frame to the chosen models.}
Fix a lawful pinned core $T$, an analytic record $R$, the algebra
tie and the source condition $\mathrm{FrameKs}$ on its models.
For row $i$, an embedding $s$ of the small field and an embedding
$\tau$ of the host, prescribe $x_0(i+1,s,\tau)$ to be the matrix
of $R.MW(i,\sigma)$ at $\rho$ if
$s=\sigma|_K$ and $\tau=\tau_E(\rho)$. Reindex using the specified
rank equality. Whenever two triples give prescriptions at one
index, the row index is determined by $i+1$, $\rho$ is determined
by injectivity of $\tau_E$, and the first clause of
$\mathrm{FrameKs}$ identifies the two model matrices. At an index
with no prescription use the old frame. On a relevant good
embedding the second clause of $\mathrm{FrameKs}$ orders the signs
according to the signature. The model's defining Gram relation
in this order gives the signed diagonal matrix in the new
frame. This constructs a single orthonormal base frame family
from the models over the small field.

The frame change has a concrete effect on the dictionary.
At an embedding where the laws apply, denote the old and new
orthonormal base matrices by $A_0,A_1$. Keep the component of
each dictionary entry and multiply its frame on the right by
$A_0^{-1}A_1$. Both base matrices are invertible by their Gram
relation. The change carries the old reference matrix to the new
one, has inverse $A_1^{-1}A_0$, preserves the corresponding
Hermitian Gram relation, and commutes with left multiplication
by every rational mover. Thus the old mover proves the lattice
and level clauses, and its archimedean frame clause follows by
cancelling $A_0 A_0^{-1}$. This gives the actual marked dictionary
change. Transport of the analytic norm laws also uses the volume
tie: the native comparison identifies the integration measures
on the moved domains. We retain that input when using the lawful
changed core below.
% END P77R1

% BEGIN P77R2
\paragraph{P77R2. The reference tuple over the embedded field itself.}
Write $e:L\hookrightarrow\mathbb C$ for the distinguished large
field embedding. Suppose that $R$ has the source preliminary
characterization $\mathrm{PreChar}$ (level, two reductions, group
record, family, models, chain and row theta laws). Take a lawful
core with its adapted $\mathrm{TieIn}$, its ordered prime equality,
the level values $\mathrm{ValLL}$, the prime guard and the definite
frames of H77S2. For the reference tuple $\mu_0$ one obtains
a pinned comparison record $O_i$ at each row with
$O_i.c'.F=e(L)$.

Indeed $e(L)$ is a number field and contains the core field by
the algebra tie. The conjugate host elements, root of unity and
lower-field inclusion in the core coefficient pin therefore give
the coefficient record over this very field. The base tuple
character is trivial. The level rule requires ramification for each
member of $Q(\mu_0)$, so $Q(\mu_0)=\varnothing$. Use the level
cover at multiplicity one with unit auxiliary ideal, auxiliary level
one and its section $1$. Its raised and lower levels agree in the
level comparisons, and the upper auxiliary factor is one.
D74U2--D74U4 give its period charts, uniformization and dictionary
with the marked base frame. The reading of the unit section on each
such frame is one by the unit evaluation law. The tuple character
has order one. The half pin of RTG10 applies with the matching
ordered primes and guard.

Here are the remaining entries. The principal level lift of a
residue group element at the marked primes gives its action on
the finite quotient of the fixed lattice. The conjugate marked
primes are disjoint by the half-level law. The two prescribed
actions on upper and lower levels follow by reducing the lift.
Use the split representatives and stabilizers of RTG6--RTG7
for the chosen frame. At every prime away from their bad modulus
the untwisted ideal coefficient is a product of the given base
and inverse reference character values (with the source norm
normalization). The $\mathrm{ValLL}$ pin places these factors
in $e(L)$. Multiplication over the prime factors of an ideal
defines its required coefficient there. RTG8--RTG9 give the
lattice coefficient from the square law at the odd row rank.
The residue test of RTG11 at $Q=\varnothing$ is a product of
zero-one indicators and an empty character product. Its values
thus lie in $e(L)$ as well.

Finally extend the chosen base coefficient functional along
$L\simeq e(L)$, and place the original definite input at the
chosen representatives via this functional. The model evaluation
on those representatives defines its adelic theta reading, as in
RTG3 and RTG5. Define the placement predicate to mean existence
of precisely such an input. The weighted coefficient on each
upper residue class is the Fourier transform of the preceding
residue test through the cover coefficient isomorphism. Its
specified equality follows by applying that isomorphism back.
These entries give the reference record's field, model vector,
level and tuple pins, both actions, class and Fourier data, and
the cover, dictionary and automorphic reading on the same row.
No enlarged auxiliary coefficient field is needed for this tuple.
% END P77R2

% BEGIN P77R3
\paragraph{P77R3. The base comparison can precede the geometric half.}
Under the conditions of P77R2 one may produce, on each of its
specified base rows, $\mathrm{OneMapOf}(c_i,1)$ for some unit
component constants $c_i$. The preliminary characterization
includes the strong theta and chain equations. In particular the
oscillator construction of Y76O1--Y76O3 applies: its level,
family, model and point-theta inputs are fields of that characterization.

To use the component comparison before there is any geometric
half, temporarily put the coordinates of a basis of $C_i$ into
the target, using the canonical cusp forms $\mathrm{formsN}$.
This gives the temporary tie $\mathrm{tieN}$: its inverse on the
coordinate image reads back the section and all its frame seminorms.
The basis extension of the canonical evaluation has the same
conditional evaluation and conjugate clauses as in A77G3, directly
from its definition. The level identity for the temporary determinant
character follows by inserting residue units into the finite idele
at one place, as in the determinant calculation of C76U1.
To obtain a face for these data, use the group record to identify
each marked residue field with its specified prime field. The
family root at that prime defines a nontrivial additive character
by sending $1$ to the primitive root. Use the trivial character
at unmarked primes. Pull back the original two finite levels and
their characters by the residue isomorphism, with the prescribed
block permutation at the second level. The auxiliary entries of
the face can be chosen at their empty or zero values; the core
only reads its group levels. The level, two reductions, group
record, family, models, chain and theta equations now give
$\mathrm{Hk3.Core}$ for this face. Replacing the target
coordinates has changed none of the inputs to the reference pins.

On this temporary record apply the reading calculation of X76R1
and X76R2. More explicitly P74I7 gives the function identity
on every placed base input, T73C5 gives its cusp section and
X76R1 gives level equivariance. They use the core just supplied,
the temporary tie, the oscillator extension and the reference
pins. The constant $uC_a$ in X76R2 has unit modulus and the
central group element and tuple multiplier can both be one.
For a section over $e(L)$, trace-dual coefficient extraction along
$K\subset e(L)$ gives the base sections exactly as in X76R2.
Projecting their translated readings with the original good
projector yields the finite $L$-linear expression in the definition
of the phase-rational clause. Equivalently, perform this last
projection on the preliminary record: the weak target coordinates
are absent from the reading sum, the determinant character in it
comes from the level residues, and its theta projection is the
strong row equation in $\mathrm{PreChar}$. Thus the same constants
and phase one solve the reference record itself. All assertions
concern the base tuple here. This supplies the base solution used
when the final geometric coordinates are constructed.
% END P77R3
